\chapter{Prefacio}

Este libro nace de una ambición sencilla: un único curso de química,
coherente de principio a fin, que abarque desde el primer año de escuela
hasta el final del grado universitario, y que pueda leerse en cualquier
parte del mundo.

El estilo es conciso y riguroso: un curso construido a partir de
definiciones, ejemplos, proposiciones, métodos y esquemas, con un
razonamiento cuidadoso siempre que resulte accesible al nivel
considerado, seguido de ejercicios cuidadosamente elegidos. Los
resultados que se enuncian sin una deducción completa se señalan
explícitamente como admitidos. Las soluciones completas de todos los
ejercicios se recogen al final del libro. Las ecuaciones químicas están
siempre ajustadas, y cada número citado para una sustancia real procede
de una tabla de referencia.

El curso consta de cuatro libros. El primero abarca los doce años de
escuela en un solo volumen y nunca enseña dos veces la misma noción: cada
idea se enseña una sola vez, en el año en que encaja por primera vez, y un
capítulo posterior que la necesite empieza con un breve recordatorio,
\emph{Lo que ya sabes}, antes de ir más allá. Los otros tres libros son
los tres años de universidad. En el primer libro, cada año escolar es una
parte, dividida en capítulos; cada libro universitario es un año,
dividido en capítulos. Cuanto más temprano es el año, más joven es el
lector al que se dirige: más ilustraciones, más pasos detallados y
ejercicios más suaves.

\section*{Cómo leer este libro}

Los enunciados van numerados dentro de cada capítulo y se distinguen por
colores: las definiciones en azul, los teoremas en rojo, las
proposiciones y sus parientes en naranja, y los métodos ---recetas
breves para las tareas habituales--- en verde. Los ejercicios están
graduados desde ($\star$), de aplicación directa, hasta
($\star\star\star$), los problemas más difíciles. Los recuadros
enmarcados recogen lo que no forma parte del razonamiento: un
recordatorio de lo que enseñó un capítulo anterior, cómo se realiza un
procedimiento en el laboratorio, una página de historia o los peligros de
un producto.

\section*{Seguridad}

La química se aprende con sustancias reales, y algunas son peligrosas.
Los experimentos descritos en este libro se realizan en un laboratorio
escolar o universitario, por un profesor o con él, con el equipo de
protección que proporciona el laboratorio. Ninguno de ellos está pensado
para hacerse en casa.

\section*{Cómo contribuir}

Las fuentes de este libro están en un repositorio git público:
\begin{center}
\url{https://github.com/bvirrion/one-chemistry-book}
\end{center}
Toda contribución ---nuevos capítulos, correcciones, mejores
explicaciones, ejercicios adicionales--- es bienvenida; véase
\texttt{CONTRIBUTING.md} en la raíz del repositorio.
